proved equivalent pole-moment formula is also acceptable if its
constants and domain are shown to match Definition~\ref{def:mu}.

\item[Prime number theorem.]
Equation~\eqref{eq:pnt}. Prove the weighted finite-interval consequence
as Lemma~\ref{lem:primesuminterface}; do not assume a stronger
normalization bound without a bridge to the actual exponents $L_p$.

\item[Integration and elementary functions.]
Complete valued fields, uniform convergence of restricted power series,
Bochner integration/Fubini/dominated convergence in the exact domains
used above, elementary gamma integrals, change of variables on the
positive half-line, and the logarithm/arctangent identities used in the
finite checker. The zero-mass energy and determinant integral have
proofs in Section~\ref{sec:real} and remain explicit project targets.

\item[Optional Legendre interface.]
Orthogonality with squared norm $2/(2j+1)$ on $[-1,1]$; parity; and the
exterior bound $|P_j(s)|\le(|s|+\sqrt{s^2-1})^j$ for $|s|\ge1$.
These support only Appendix~\ref{app:approx}.
\end{description}

\subsection{Representation and statement-level hazards}
\begin{enumerate}[leftmargin=*,label=\arabic*.]
\item Work with the integer parameter $n$ and set $K=40n$, $N=3n$,
$h=37n$. This avoids accidental natural-number division or unproved
divisibility. Represent $Q_{K,M}$ on an admissible parameter subtype, or
supply harmless values outside the admissible domain and prove the
bridge on that domain.
\item Define the rational functional using its unique polynomial and
simple-pole decomposition. The domain is a vector space; arbitrary
products may introduce double poles and leave the domain. Extend
scalars explicitly before applying it over $\Qp$.
\item Use a restricted sequence space plus a finite residue vector if
that avoids a general Tate-algebra library. Preserve the bounded
analytic remainder and the near/far distinction. The target of a
convergent local evaluation has bounded polynomial degree; do not
silently use completeness of all of $\Qp[X]$ in the Gauss norm.
\item State all arithmetic bounds for the \emph{Gauss valuation of
coefficients}, including the constant coefficient, not merely for the
specialized value at $\xi$. Keep the small primes $2,3,5$ in the
universal argument and retain the factor $v_p(24)$.
\item Keep the signs of $L_p$ and use integer powers in \eqref{eq:mdef}.
The inner row weights may be positive. Only the rank-correction lemma
requires all its weights to be nonpositive. Its rank hypothesis is rank
over $\Qp$, not rank of a reduction modulo $p$.
\item Keep $M$ fixed before passing to $K\to\infty$. Distinguish the
explicit admissibility threshold from the eventual asymptotic threshold.
A numerical check at $n=200000$ is not a proof of eventual decay.
\item In standard Lean real analysis, \texttt{Real.log 0} is a totalized
value, not $-\infty$ (see \cite{MathlibLog}). The logarithmic-energy interface must therefore
include nullity of the diagonal as well as integrability of a real-valued
representative, or use an extended-valued singular kernel and prove that
nullity. For configuration bounds first assume the points are pairwise
distinct; treat collisions separately by the zero Vandermonde product.
Do not translate the displayed $-\infty$ convention into a false real
inequality involving a totalized logarithm.
\item Before every inverse, logarithm, and determinant quotient, retain
the needed nonzero or positivity proof. The source weight is defined
for $y>0$; a value at $y=0$ must be assigned separately if a globally
continuous extension is desired.
\item Validate every certificate branch, including square-root domains,
reciprocal nonvanishing, partition coverage, and the finite set of
excluded integral endpoints. Validating only final floating-point values
is insufficient. In the outer certificate retain the entire interval
$(1,2\lambda)$ and the piecewise definition \eqref{eq:douter}.
\end{enumerate}

\subsection{Recommended build stages and completion criteria}
\begin{center}\small
\begin{tabular}{>{\raggedright\arraybackslash}p{0.08\linewidth}>{\raggedright\arraybackslash}p{0.27\linewidth}>{\raggedright\arraybackslash}p{0.54\linewidth}}
\toprule
Stage & Module group & Deliverable\\
\midrule
1 & Definitions and endpoint & Actual determinant definitions,
real-series target, rational-evaluation lemma, and a visibly conditional
end-to-end theorem.\\
2 & Exact arithmetic & Rational constants, parameter identities,
Appendix B branch certificates, and a sound interval checker.\\
3 & Completed local functional & Bounded division, principal-part
compatibility, reflection/difference, and the full distribution identity.\\
4 & Local determinant bounds & Small-prime bound, both unimodular
bases, all weight comparisons, field-rank correction, integrality.\\
5 & Prime normalization & Uniform finite-count estimates, PNT interface,
weighted prime sums, periodic-tail bound, fixed-cutoff growth.\\
6 & Positive and real analysis & Moment representation, Gaussian
zero-mass energy, regularization, confinement, determinant integral.\\
7 & Main endpoint & Instantiate every input; prove integer polynomial
decay and the unconditional irrationality theorem; audit axioms.\\
8 & Optional consequences & Relative norm, exponent $260$, height,
primitive normalization.\\
\bottomrule
\end{tabular}
\end{center}
The stages may proceed in parallel where dependencies allow. The local
functional and numerical certificates deserve early attention because
they resolve the most compressed parts of the source argument.

The files \path{formalization/declarations.csv} and
\path{formalization/dependency_graph.md} give the stable target list and
its dependencies. \path{AGENTS.md} supplies operational instructions.
No proposed theorem may be weakened silently to make a build pass.
When an obstruction occurs, record the exact failed statement and a
minimal counterexample or missing lemma; distinguish a mathematical
counterexample from missing library infrastructure. Keep public status
conditional until the endpoint has no project-owned assumptions.

\subsection{Verification performed for this edition}
The finite checker was executed with ordinary Python integer arithmetic
and 192-bit outward-rounded intervals. It accepted all sixteen measure
components, the minimizing-point bracket and its sign conditions, all
684 potential cells, the energy and scalar enclosures, the optional
$\kappa$ enclosure, the 143 symbolic inner branches, all seventeen unit
integrals, the eleven outer branches, and both exact final margins.
It also checks the elementary rational tail calculation at $20$ and the
two approximation-constant identities.

The provided audits additionally report small-$K$ determinant
calculations and truncated $p$-adic distribution tests. Those reports
are preserved as provenance in \path{REPAIRS.md}; they have not been
relabelled as computations performed in this edition, and none is used
as a proof premise. This document does not assert that all possible
mathematical defects have been excluded. It makes the revised proof
obligations explicit and supplies reproducible finite certificates for
the computations it actually checks.

\begin{thebibliography}{9}
\bibitem{Fauzan}
A. Fauzan, \emph{$\zeta(5)$ is irrational}, supplied 32-page manuscript,
17 September 2026. The source file and its SHA-256 digest are identified
in Section~\ref{sec:overview}. Original author contact in the source:
\texttt{aabir.fauzan@aalto.fi}.

\bibitem{DLMF}
National Institute of Standards and Technology,
\emph{Digital Library of Mathematical Functions}, Chapters 18, 24, 25,
and 27. In particular (18.10.5), (24.4.1), (24.4.3), (24.4.18),
\S24.10(i), (25.6.2), (25.11.29), and \S27.12.
\url{https://dlmf.nist.gov/}.

\bibitem{Saff}
E. B. Saff, \emph{Logarithmic potential theory with applications to
approximation theory}, Surveys in Approximation Theory \textbf{5}
(2010), 165--200, arXiv:1010.3760.

\bibitem{Forrester}
P. J. Forrester, \emph{Meet Andr\'eief, Bordeaux 1886, and Andreev,
Kharkov 1882--1883}, Random Matrices: Theory and Applications
\textbf{8} (2019), no. 2, 1930001, arXiv:1806.10411.
\bibitem{LeanValidation}
The Lean reference manual, \emph{Axioms} and \emph{Validating a Lean Proof},
consulted 23 September 2026.
\url{https://lean-lang.org/doc/reference/latest/Axioms/};
\url{https://lean-lang.org/doc/reference/latest/ValidatingProofs/}.

\bibitem{MathlibLog}
The Mathlib documentation, \texttt{Mathlib.Analysis.SpecialFunctions.Log.Basic},
in particular \texttt{Real.log\_zero}, consulted 23 September 2026.
\url{https://leanprover-community.github.io/mathlib4_docs/Mathlib/Analysis/SpecialFunctions/Log/Basic.html}.
\end{thebibliography}
\end{document}
